\pdfoutput=1
\documentclass[11pt]{amsart}
\usepackage[T1]{fontenc}
\usepackage[utf8]{inputenc}
\usepackage{lmodern}
\usepackage[margin=1.15in]{geometry}
\usepackage{amsmath,amssymb,amsthm,mathtools}
\usepackage{booktabs}
\usepackage[section]{placeins}
\usepackage{microtype}
\usepackage{enumitem}
\usepackage{xcolor}
\usepackage[hidelinks]{hyperref}
\usepackage{xurl}
\usepackage{orcidlink}
\usepackage[capitalise,nameinlink]{cleveref}

\newtheorem{theorem}{Theorem}[section]
\newtheorem{lemma}[theorem]{Lemma}
\newtheorem{corollary}[theorem]{Corollary}
\newtheorem{proposition}[theorem]{Proposition}
\theoremstyle{remark}
\newtheorem{remark}[theorem]{Remark}

\newcommand{\code}[1]{\texttt{#1}}
\newcommand{\sv}{\ell_2}

\title[Power-of-two cycles in censuses of cubic graphs]{Cycles of length 4, 8, 16 and 32\\ in censuses of cubic graphs}

\author[A. Marszałek]{Andrzej Marszałek\,\orcidlink{0009-0006-6148-2965}}
\address{Independent Researcher}
\urladdr{https://github.com/hipotures/conjectures}

\subjclass[2020]{05C38, 05C25, 05C35, 68V05}
\keywords{Erdős–Gyárfás conjecture, cubic graphs, vertex-transitive graphs, arc-transitive graphs, graph censuses, girth}

\date{27 November 2026}

\begin{document}

\begin{abstract}
The Erdős–Gyárfás conjecture states that every graph of minimum degree at least
three contains a cycle whose length is a power of two. We check it, with
explicit and independently verified witnesses, on three published lists of
cubic graphs. All $1{,}436{,}474$ connected cubic graphs of girth at least $9$
on $58$ to $64$ vertices contain a $16$-cycle. All $111{,}360$ graphs of the
Potočnik--Spiga--Verret census of cubic vertex-transitive graphs on at most
$1280$ vertices, and all $3815$ graphs of Conder's list of cubic symmetric
graphs on at most $10{,}000$ vertices, contain a cycle of length $4$, $8$, $16$
or $32$. We determine the shortest power-of-two cycle length of every graph.
It is $32$ for exactly $100$ vertex-transitive graphs, the smallest of which has
$180$ vertices, and for exactly $689$ symmetric graphs, the smallest of which has
$1050$ vertices. Among symmetric graphs the share that needs a $32$-cycle grows
from about $1\%$ below $1280$ vertices to about $25\%$ between $5000$ and
$10{,}000$ vertices.
\end{abstract}

\maketitle

%======================================================================
\section{Introduction}

\begin{quote}
\textbf{Conjecture} (Erdős and Gyárfás, 1995; see~\cite{Erdos1997}).
Every finite simple graph of minimum degree at least three contains a cycle
whose length is a power of two.
\end{quote}

The conjecture is open even for cubic graphs. Exhaustive searches exclude small
counterexamples: a counterexample has more than $40$ vertices, and more than
$48$ if it is cubic~\cite{DucoffeDumitru2026}; a cubic bipartite one has at
least $86$ vertices~\cite{Marszalek2026b}. Exoo~\cite{Exoo2014} found a cubic
graph on $78$ vertices without cycles of length $4$, $8$ and $16$, and Ripatti's
search~\cite{Ripatti2017} indicates that no smaller cubic graph has this
property. Among structured classes, Burkhart~\cite{Burkhart2026} proved the
conjecture for Cayley graphs of groups of odd order and of nilpotent groups.
Cubic Cayley graphs of other groups, and cubic vertex-transitive graphs that are
not Cayley graphs, are not covered. Gebendorfer~\cite{Gebendorfer2026} studied
the cubic vertex-transitive bipartite graphs of girth $12$ without a $16$-cycle.

Graphs of large girth are natural places to look for counterexamples, since
they have no short cycles at all, and many of the smallest known cubic graphs
of a given girth are vertex-transitive or arc-transitive. Such graphs are also
the usual starting points for constructions that avoid short power-of-two
cycles, such as voltage lifts and vertex replacements~\cite{Exoo2014,Marszalek2026}.
Complete lists of these graphs exist up to certain orders, and we check all of
them.

All graphs are finite and simple. For a graph $G$ let $\sv(G)$ be the least power of two $2^k\ge 4$ such that $G$
has a cycle of length $2^k$ (and $\sv(G)=\infty$ if there is none).

\begin{theorem}\label{thm:girth9}
Every connected cubic graph of girth at least $9$ on at most $64$ vertices
contains a $16$-cycle.
\end{theorem}

\begin{corollary}\label{cor:girth8}
Every cubic graph of girth at least $8$ on at most $64$ vertices contains a
cycle of length $8$ or $16$.
\end{corollary}

\begin{proof}
Let $G$ be such a graph without an $8$-cycle. Then $G$ has girth at least $9$,
so each component has at least $58$ vertices~\cite{BrinkmannMcKaySaager1995},
and $G$ is connected. \Cref{thm:girth9} applies.
\end{proof}

\begin{theorem}\label{thm:vt}
Every graph in the Potočnik--Spiga--Verret census of connected cubic
vertex-transitive graphs on at most $1280$ vertices~\cite{PSV2013} contains a
cycle of length $4$, $8$, $16$ or $32$. Exactly $100$ of the $111{,}360$ graphs
have $\sv(G)=32$; the smallest has $180$ vertices.
\end{theorem}

\begin{theorem}\label{thm:sym}
Every graph in Conder's list of connected cubic symmetric (arc-transitive)
graphs on at most $10{,}000$ vertices~\cite{ConderList} contains a cycle of
length $4$, $8$, $16$ or $32$. Exactly $689$ of the $3815$ graphs have
$\sv(G)=32$; the smallest has $1050$ vertices.
\end{theorem}

\Cref{thm:girth9} assumes that the House of Graphs lists of cubic graphs of
girth at least $9$~\cite{HoG} are complete, and \cref{thm:vt,thm:sym} are
statements about the two censuses as published; we do not regenerate them. Since
the censuses are complete~\cite{PSV2013,ConderList}, \cref{thm:vt,thm:sym} say
that every connected cubic vertex-transitive graph on at most $1280$ vertices,
and every connected cubic symmetric graph on at most $10{,}000$ vertices, has a
cycle of length $4$, $8$, $16$ or $32$. If Ripatti's result~\cite{Ripatti2017}
is correct, it implies \cref{thm:girth9} (but not \cref{thm:vt,thm:sym}, whose
graphs are larger); our proof is independent of it.

\Cref{tab:summary} summarises the three computations.

\begin{table}[ht]
\centering
\caption{The three censuses and the distribution of $\sv(G)$.}
\label{tab:summary}
\begin{tabular}{@{}lrrrrrr@{}}
\toprule
census & orders & graphs & $\sv=4$ & $\sv=8$ & $\sv=16$ & $\sv=32$\\
\midrule
cubic, girth $\ge 9$~\cite{HoG} & $58$--$64$ & $1{,}436{,}474$ & -- & -- & $1{,}436{,}474$ & --\\
vertex-transitive~\cite{PSV2013} & $4$--$1280$ & $111{,}360$ & $5756$ & $16{,}624$ & $88{,}880$ & $100$\\
symmetric~\cite{ConderList} & $4$--$10{,}000$ & $3815$ & $3$ & $111$ & $3012$ & $689$\\
\bottomrule
\end{tabular}
\end{table}

\subsection*{Nature of the proof}
The proofs are computer-assisted. A search program finds, for each graph, a
cycle of length $\sv(G)$ and reports it. A second program, written separately,
reads the original census files with its own parsers, checks every reported
cycle edge by edge, and confirms with a different search algorithm that no
shorter power-of-two cycle exists. The positive part of each statement (a cycle
exists) therefore rests only on the census files and on a check that a given
vertex sequence is a cycle. It does not depend on vertex-transitivity
either. The minimality of $\sv(G)$ is not needed for these existence statements
(or for \cref{thm:girth9}), but it is needed for the exact counts and the
smallest orders in \cref{thm:vt,thm:sym}. For the vertex-transitive and
symmetric censuses it rests on the two searches agreeing and on
vertex-transitivity, through \cref{lem:vt}.

%======================================================================
\section{The censuses}\label{sec:censuses}

All input files are fixed by their SHA-256 hashes in the repository
(\cref{sec:repro}).

\emph{Cubic graphs of girth at least $9$.} House of Graphs~\cite{HoG} provides
the connected cubic graphs of girth at least $9$ on $58$, $60$, $62$ and $64$
vertices: $18$, $474$, $27{,}169$ and $1{,}408{,}813$ graphs, in graph6 format.
There is none on fewer than $58$ vertices, and the $18$ graphs on $58$
vertices are the $(3,9)$-cages~\cite{BrinkmannMcKaySaager1995}.

\emph{Cubic vertex-transitive graphs.} Potočnik, Spiga and Verret~\cite{PSV2013}
determined all connected cubic vertex-transitive graphs on at most $1280$
vertices. We use the database release~\cite{PSVdata}, whose file
\code{main\_Graph.txt} contains $111{,}705$ rows. The $111{,}360$ rows with a
census index are the census; the other $345$ rows carry no census index and
no adjacency lists, and we do not use them. We write $\mathrm{CVT}(n,i)$ for
the graph of order $n$ with census index $i$. Each graph is read from its
adjacency lists (the column \code{raw\_magma\_code}); the census also records
the girth of every graph, which we compare with our own computation.

\emph{Cubic symmetric graphs.} Conder~\cite{ConderList} lists all connected
cubic arc-transitive graphs on at most $10{,}000$ vertices, extending the list
of Conder and Dobcsányi~\cite{ConderDobcsanyi2002}. The file gives each graph as
a set of edges, with a header line stating its name (such as
$\mathrm{C1050.1}$), girth and other invariants.

%======================================================================
\section{Method}\label{sec:method}

\subsection{Search}
Let $G$ be a cubic graph and $v_0$ a vertex. For $h\ge 2$ a cycle of length
$2h$ through $v_0$ is the union of two paths of length $h$ from $v_0$ to the
vertex $w$ opposite $v_0$ on the cycle, and the two paths have disjoint sets of
inner vertices. Conversely:

\begin{lemma}\label{lem:mitm}
Let $P=v_0x_1\cdots x_{h-1}w$ and $Q=v_0y_1\cdots y_{h-1}w$ be paths with the same
ends and $\{x_1,\dots,x_{h-1}\}\cap\{y_1,\dots,y_{h-1}\}=\emptyset$. Then
$v_0x_1\cdots x_{h-1}wy_{h-1}\cdots y_1$ is a cycle of length $2h$.
\end{lemma}

\begin{proof}
The $2h$ listed vertices are distinct: $v_0$ and $w$ are ends of both paths and
so occur in neither interior, and the interiors are disjoint. Consecutive
vertices, including $y_1$ and $v_0$, are adjacent.
\end{proof}

The search program \code{census\_scan} therefore lists all paths of length $h$
from $v_0$ (at most $3\cdot 2^{h-1}$ in a cubic graph), sorts them by end
vertex, and looks for two paths with the same end and disjoint interiors. This
finds a $2h$-cycle through $v_0$ whenever one exists. The program tries
$h=2,4,8,16$ in turn and stops at the first $h$ for which a cycle exists.

\begin{lemma}\label{lem:vt}
If $G$ is vertex-transitive and has a cycle of length $L$, then every vertex
of $G$ lies on a cycle of length $L$.
\end{lemma}

\begin{proof}
An automorphism maps a vertex of the cycle to any given vertex, and maps the
cycle to a cycle of the same length.
\end{proof}

For the vertex-transitive and symmetric censuses the program therefore searches
through $v_0=0$ only, and $\sv(G)$ is the least $L$ for which a cycle through
$0$ is found. For the girth-$9$ graphs it searches through every vertex. In
each case it reports $\sv(G)$, the girth, and one cycle of length $\sv(G)$ as
a list of vertices.

\subsection{Checking}
The checker \code{check\_witnesses.py} does not use the converted files that the search
program reads. It reads the census files with its own parsers (a character scanner for
the adjacency lists, a line-based reader for Conder's file, and its own graph6
decoder) and, for every graph:
\begin{enumerate}[label=(\roman*)]
\item checks that the graph is simple, cubic and connected;
\item checks that the reported cycle has $\sv(G)$ distinct vertices and that
consecutive vertices, including the last and the first, are adjacent;
\item computes the girth and compares it with the search program and, for the
vertex-transitive and symmetric censuses, with the girth stored in the census;
\item confirms minimality. For the girth-$9$ graphs the girth computed here
(from every vertex) is at least $9$, which excludes $4$- and $8$-cycles. For the
other two censuses, a depth-first search over paths $0=x_0,x_1,\dots$ pruned by
$\operatorname{dist}(x_d,0)\le L-d$ finds no cycle of length $L$ through $0$ for
each power of two $L<\sv(G)$, and finds one for $L=\sv(G)$.
\end{enumerate}
It also checks that every graph has exactly one result line and that the number
of graphs agrees with \cref{sec:censuses}.

The girth is computed as follows: for a vertex $s$, a breadth-first search
labels each vertex by the neighbour of $s$ through which it is first reached;
the least value of $d(u)+d(w)+1$ over edges $uw$ not incident to $s$ with
differently labelled ends is the length of a shortest cycle through $s$. Indeed,
both tree paths together with $uw$ form a cycle; and if $C=sv_1\cdots v_{k-1}$
is a shortest cycle through $s$, then $v_1$ and $v_{k-1}$ have different labels,
so the label changes on some edge $v_iv_{i+1}$ with $1\le i<k-1$, and every such
edge satisfies $d(v_i)+d(v_{i+1})+1\le |C|$. Vertex-transitivity is taken
from the census and not checked.

%======================================================================
\section{Results}\label{sec:results}

\Cref{thm:girth9}: every one of the $1{,}436{,}474$ graphs has $\sv(G)=16$.

\Cref{tab:vt,tab:sym} give $\sv(G)$ by girth for the vertex-transitive and the symmetric
census.
Rows with girth $4$, $8$ or $16$ are immediate, since then $\sv(G)$ equals the
girth; graphs of girth at least $17$ have $\sv(G)\ge 32$ for the same reason.

\begin{table}[ht]
\centering
\caption{Cubic vertex-transitive graphs on at most $1280$ vertices: number of
graphs by girth $g$ and shortest power-of-two cycle length $\sv$.}
\label{tab:vt}
\begin{tabular}{@{}rrrrrr@{}}
\toprule
$g$ & $\sv=4$ & $\sv=8$ & $\sv=16$ & $\sv=32$ & total\\
\midrule
3 & 2 & 3 & 102 & 53 & 160\\
4 & 5754 & & & & 5754\\
5 & & 14 & 83 & 3 & 100\\
6 & & 2957 & 55{,}692 & 25 & 58{,}674\\
7 & & 121 & 70 & 1 & 192\\
8 & & 13{,}529 & & & 13{,}529\\
9 & & & 218 & 1 & 219\\
10 & & & 25{,}805 & 1 & 25{,}806\\
11 & & & 80 & & 80\\
12 & & & 5407 & 16 & 5423\\
13 & & & 37 & & 37\\
14 & & & 1365 & & 1365\\
15 & & & 12 & & 12\\
16 & & & 9 & & 9\\
\addlinespace
total & 5756 & 16{,}624 & 88{,}880 & 100 & 111{,}360\\
\bottomrule
\end{tabular}
\end{table}

\begin{table}[ht]
\centering
\caption{Cubic symmetric graphs on at most $10{,}000$ vertices: number of
graphs by girth $g$ and shortest power-of-two cycle length $\sv$.}
\label{tab:sym}
\begin{tabular}{@{}rrrrrr@{}}
\toprule
$g$ & $\sv=4$ & $\sv=8$ & $\sv=16$ & $\sv=32$ & total\\
\midrule
3 & 1 & & & & 1\\
4 & 2 & & & & 2\\
5 & & 2 & & & 2\\
6 & & 7 & 1559 & & 1566\\
7 & & 2 & 8 & 9 & 19\\
8 & & 100 & & & 100\\
9 & & & 40 & & 40\\
10 & & & 290 & 21 & 311\\
11 & & & 7 & 4 & 11\\
12 & & & 571 & 310 & 881\\
13 & & & 2 & 1 & 3\\
14 & & & 248 & 125 & 373\\
15 & & & 22 & 14 & 36\\
16 & & & 265 & & 265\\
17--20 & & & & 205 & 205\\
\addlinespace
total & 3 & 111 & 3012 & 689 & 3815\\
\bottomrule
\end{tabular}
\end{table}

\emph{Smallest examples.} The smallest vertex-transitive graph with $\sv(G)=32$
is $\mathrm{CVT}(180,50)$: it has girth $3$, is a Cayley graph and is not
arc-transitive. The $53$ graphs of girth $3$ with $\sv(G)=32$ are the most
numerous group; the smallest graph with $\sv(G)=32$ and girth at least $4$ is
$\mathrm{CVT}(504,257)$, of girth $6$. The smallest symmetric graph with
$\sv(G)=32$ is $\mathrm{C1050.1}$, bipartite of girth $12$. Six symmetric graphs
on at most $1280$ vertices have $\sv(G)=32$ ($\mathrm{C1050.1}$,
$\mathrm{C1134.3}$, $\mathrm{C1134.4}$, $\mathrm{C1134.5}$, $\mathrm{C1140.4}$,
$\mathrm{C1248.1}$), and they are exactly the arc-transitive graphs among the
$100$ of \cref{thm:vt}. The file \code{data/shortest32.tsv} lists all $789$
census entries with $\sv(G)=32$ together with a $32$-cycle; they are $783$
distinct graphs, since the six graphs above appear in both censuses.

\emph{Growth with the order.} Among symmetric graphs, the share with
$\sv(G)=32$ is $1.2\%$ ($6$ of $482$) up to $1280$ vertices, $6.4\%$ ($31$ of
$487$) from $1281$ to $2560$, $18.8\%$ ($177$ of $942$) from $2561$ to $5000$,
and $24.9\%$ ($475$ of $1904$) from $5001$ to $10{,}000$ vertices. Part of
this growth comes from girth alone: $205$ of the $689$ graphs have girth $17$ to
$20$, where $\sv(G)\ge 32$ is forced, and all of them have more than $2560$
vertices. Among symmetric graphs of girth at most $16$ the share is $15.8\%$
($144$ of $909$) from $2561$ to $5000$ and $17.5\%$ ($303$ of $1732$) from $5001$
to $10{,}000$ vertices. Among
vertex-transitive graphs it stays below $0.11\%$ in each quarter of the range
up to $1280$ vertices. No graph in either census needs a $64$-cycle.

\emph{Girth $12$.} The $16$ graphs of girth $12$ in \cref{tab:vt} with
$\sv(G)=32$ are all bipartite, and no non-bipartite vertex-transitive graph of
girth $12$ on at most $1280$ vertices lacks a $16$-cycle. This agrees with
Gebendorfer's count of exactly sixteen cubic vertex-transitive bipartite graphs
of girth twelve without a $16$-cycle in the same census, each with a
$32$-cycle~\cite{Gebendorfer2026}.

%======================================================================
\section{Validation}\label{sec:validation}

Besides the checker of \cref{sec:method}, we ran the following checks.
\begin{enumerate}[label=(\alph*)]
\item \emph{Census metadata.} The girth computed by both programs agrees with
the girth stored in the census for all $111{,}360$ vertex-transitive and all
$3815$ symmetric graphs.
\item \emph{Negative controls.} The script \code{test\_checker.py} feeds the
checker corrupted results for the Möbius--Kantor graph and the Heawood graph: a
claimed $16$-cycle where an $8$-cycle exists, a cycle of the wrong length, a
non-adjacent step, a repeated vertex, a wrong girth (reported or in the
census), a mismatched graph name, a missing cycle, and a graph of girth $6$
offered as a girth-$9$ graph. The checker rejects each of them and accepts the
two correct results.
\item \emph{Prior counts.} The census sizes agree with the published
ones~\cite{PSV2013,ConderList,HoG}, and the girth-$12$ count agrees
with~\cite{Gebendorfer2026}.
\end{enumerate}

%======================================================================
\section{Reproducibility}\label{sec:repro}

Everything is in the repository \url{https://github.com/hipotures/conjectures},
in the folder \code{results/}\allowbreak\code{cubic-census-power-of-two-cycles},
and is in the public domain (CC0). The script \code{run\_census.py} downloads the six input files (about
$2.6$\,GB, mostly the vertex-transitive census) and checks their SHA-256 hashes
against \code{data/sources.json}. It then compiles and runs the search program
\code{src/census\_scan.cpp}, runs the checker \code{check\_witnesses.py}, and
writes the files \code{data/summary.json} and \code{data/shortest32.tsv}. Once the
files are downloaded, a run takes a little over three minutes on a $16$-core
machine: about a minute to convert the files, half a minute for the searches
and a minute and a half for the checks.

\section*{Use of AI tools}
This work was carried out with substantial assistance from AI systems, mainly
Gemini~4~Neon (Google) and Claude Opus~5.5 (Anthropic), used through the
Antigravity and Claude Code environments. They were used to write and run the
search and checking programs and to draft the text. AI systems are not authors
of this paper. The author takes responsibility for the content and for the
attribution of prior work.

\section*{Acknowledgements}
This work relies entirely on the censuses of Primož Potočnik, Pablo Spiga and
Gabriel Verret, of Marston Conder, and of House of Graphs, and on their public
release.

\begin{thebibliography}{99}

\bibitem{PSVdata}
J.~Azarija, K.~Berčič, F.~Koprivec, T.~Marc, P.~Potočnik, P.~Rotar, P.~Spiga,
G.~Verret and J.~Vidali, \emph{Cubic vertex-transitive graphs on up to 1280
vertices}, dataset, Zenodo, 2022, \url{https://doi.org/10.5281/zenodo.6576526}.

\bibitem{BrinkmannMcKaySaager1995}
G.~Brinkmann, B.~D.~McKay and C.~Saager, \emph{The smallest cubic graphs of
girth nine}, Combin. Probab. Comput. \textbf{4} (1995), 317--329,
\url{https://doi.org/10.1017/S0963548300001693}.

\bibitem{Burkhart2026}
M.~C.~Burkhart, \emph{Odd-ordered Cayley graphs satisfy Erdős--Gyárfás},
arXiv:2611.05060, 2026.

\bibitem{ConderList}
M.~Conder, \emph{Trivalent (cubic) symmetric graphs on up to 10000 vertices},
list of edge sets, 2024,
\url{https://www.math.auckland.ac.nz/~conder/AllSymmCubic-OrderUpTo10000-ByEdgeSets.txt.gz}.

\bibitem{ConderDobcsanyi2002}
M.~Conder and P.~Dobcsányi, \emph{Trivalent symmetric graphs on up to 768
vertices}, J. Combin. Math. Combin. Comput. \textbf{40} (2002), 41--63.

\bibitem{HoG}
K.~Coolsaet, S.~D'hondt and J.~Goedgebeur, \emph{House of Graphs 2.0: a
database of interesting graphs and more}, Discrete Appl. Math. \textbf{325}
(2023), 97--107, \url{https://doi.org/10.1016/j.dam.2022.10.013}. Lists:
\url{https://houseofgraphs.org/meta-directory/cubic}.

\bibitem{DucoffeDumitru2026}
G.~Ducoffe and B.~Dumitru, \emph{Towards a more structured search for
Erdős--Gyárfás counter-examples}, arXiv:2609.28594, 2026.

\bibitem{Erdos1997}
P.~Erdős, \emph{Some old and new problems in various branches of combinatorics},
Discrete Math. \textbf{165/166} (1997), 227--231.

\bibitem{Exoo2014}
G.~Exoo, \emph{Three graphs and the Erdős--Gyárfás conjecture},
arXiv:1403.5636, 2014.

\bibitem{Gebendorfer2026}
J.~J.~Gebendorfer, \emph{The Erdős--Gyárfás conjecture for cubic
vertex-transitive graphs of girth twelve without sixteen-cycles: flat covers,
chains, and the dodecagon-factor frontier}, preprint, Zenodo, 2026,
\url{https://doi.org/10.5281/zenodo.21295085}.

\bibitem{Marszalek2026}
A.~Marszałek, \emph{Small graphs without cycles of length a power of two: new
bounds with kernel-checked proofs}, preprint, 2026,
\url{https://doi.org/10.5281/zenodo.24588313}.

\bibitem{Marszalek2026b}
A.~Marszałek, \emph{Every cubic bipartite graph on at most 84 vertices has a
cycle of length 4, 8 or 16}, preprint, 2026,
\url{https://doi.org/10.5281/zenodo.24599605}.

\bibitem{PSV2013}
P.~Potočnik, P.~Spiga and G.~Verret, \emph{Cubic vertex-transitive graphs on up
to 1280 vertices}, J. Symbolic Comput. \textbf{50} (2013), 465--477,
\url{https://doi.org/10.1016/j.jsc.2012.09.002}.

\bibitem{Ripatti2017}
A.~V.~Ripatti, \emph{Faster approach for exhaustive search over graphs with
certain properties}, abstract, BAIKAL 2017 conference, 2017.

\end{thebibliography}

\end{document}
